\section{A separate adequate-resource specialization}\label{app:adequate_resources}
We next study a different hardware policy in which the complete assignment set has an exact
energy projection. There are $n\ge1$ mandatory depot-to-depot services
$A_i$ occupying $[0,1]$ and $n$ mandatory services $B_j$ occupying $[2,3]$.
Their consumptions are $a_i$ and $b_j$, with
\begin{equation}\label{eq:batchassumptions}
 0<a_i\le20,\qquad0<b_j\le20,\qquad a_i+b_j\le30
 \quad\text{for every }i,j.
\end{equation}
Every A-to-B assignment is allowed, without movement time, energy or cost.
At most $2n$ identical buses may be used. Each has capacity 20 kWh,
starts full and returns full by hour 4. Its operating fee is strictly
positive, $f>0$.

Early charging in $[1,2]$ uses $n$ connectors of 10 kW each and an
aggregate grid limit of $10n$ kW. Terminal charging in $[3,4]$ uses $n$
connectors of 30 kW each and an aggregate grid limit of $30n$ kW. Per-bus
acceptance matches those connector powers. There is no other charging,
reserve, loss, taper or switching time. The grid limits are part of the
assumption: connector counts alone would not establish the simultaneous
allocation used below.

Since the services in each batch are simultaneous and have positive duration,
a bus can cover at most one A and one B. Thus a complete assignment is a
partial bipartite matching $M$. A matched pair uses one bus; unmatched
services each use a separate bus. All crossmatchings and all cardinalities
$m=0,\ldots,n$ are retained.

Define
\begin{align}
 S&=\sum_i a_i+\sum_j b_j,&u_i&=\min\{a_i,10\},&U&=\sum_i u_i,\\
 r_{ij}&=[a_i+b_j-20]_+,&
 r_m&=\min_{|M|=m}\sum_{(i,j)\in M}r_{ij},&c_m&=f(2n-m),
\end{align}
where $r_0=0$. The edge cost $r_{ij}$ is the minimum early purchase
needed by a bus covering $A_i$ then $B_j$. The thresholds $r_m$ are
minimum-cost cardinality-matching values, not samples from a plan bank.

\begin{theorem}[Exact heterogeneous fleet projection]\label{thm:batchprojection}
Under \eqref{eq:batchassumptions} and the stated resources, including $f>0$,
the complete physical projection is exactly
\begin{equation}\label{eq:batchprojection}
 \bigcup_{m=0}^n\{(E,S-E,c_m):r_m\le E\le U\}.
\end{equation}
The thresholds are nondecreasing and discretely convex:
\begin{equation}\label{eq:discreteconvex}
 0\le r_1-r_0\le r_2-r_1\le\cdots\le r_n-r_{n-1}.
\end{equation}
Let $m_0=\max\{m:r_m=0\}$. The lowest operating-cost boundary of the
complete-plan hull is the affine interpolant through
$(r_m,c_m)$ for $m=m_0,\ldots,n$, followed by the horizontal segment
$c_n$ on $[r_n,U]$.
\end{theorem}

The physical part of the theorem is the battery and connector argument.
A bus that has served $A_i$ can buy any $x_i\le u_i$ early; pairing it
with $B_j$ adds the lower bound $x_i\ge r_{ij}$. There are exactly $n$
A buses, so their individual charge intervals can be allocated simultaneously.
Terminal replenishment fits by grouping each matched bus alone and each
unmatched A/B pair together on one connector. The discrete-convexity step
uses standard bipartite min-cost-flow integrality and convex LP value
functions. It is not a new matching theorem. The complete proof is in
Appendix~\ref{app:heterogeneous}.

Use the same scaled supply family as the later replication calculation:
\begin{equation}\label{eq:batchsupply}
 F_n(E,L)=4E+\frac{E^2+L^2}{10n},\qquad
 G(E)=F_n(E,S-E),\qquad E_0=\frac{S-20n}{2}.
\end{equation}
Strict convexity makes the hull minimizing early load $E^*$ unique.

\begin{theorem}[A complete support criterion for the two-batch class]
\label{thm:batchsupport}
For strictly positive fee $f>0$ under Theorem~\ref{thm:batchprojection},
\begin{equation}\label{eq:batchzero}
 D=\CH\quad\Longleftrightarrow\quad
 E^*\in\{r_m:m=m_0,\ldots,n\}\cup[r_n,U].
\end{equation}
The gap is strictly positive exactly when $E^*$ lies strictly inside an
adjacent threshold interval $(r_{j-1},r_j)$, $j>m_0$.
If $d_j=r_j-r_{j-1}>0$, an intermediate threshold $r_m$ minimizes the
hull objective exactly when
\begin{equation}\label{eq:batchvertex}
 \frac{f}{d_{m+1}}\le G'(r_m)\le\frac{f}{d_m},
 \qquad m_0<m<n.
\end{equation}
\end{theorem}

The criterion includes weak support at ties. A threshold is physical, but
an interior point between consecutive thresholds interpolates operating
costs that no fleet at that early load can attain. Collinear adjacent
segments do not erase the intervening physical threshold. If $f=0$,
the conclusion about strict interiors would fail: the zero-matching
branch realizes every early load at zero operating cost, so the gap is
always zero. Strict positivity of the fee is therefore substantive.

The values themselves require only matching thresholds and scalar convex
minimizations. With $\clamp(v;\ell,u)=\min\{u,\max\{\ell,v\}\}$,
\begin{equation}\label{eq:batchD}
 D=\min_{0\le m\le n}\{c_m+G(\clamp(E_0;r_m,U))\}.
\end{equation}
For each sloping hull edge $j>m_0$, evaluate its objective at
\begin{equation}\label{eq:batchh}
 h_j=\clamp\left(E_0+\frac{5nf}{2d_j};r_{j-1},r_j\right).
\end{equation}
Compare these values with the minimum on the final horizontal segment.
Their minimum is $\CH$. If $m_0=n$, that horizontal segment alone
remains and $D=\CH$.


The separate eight-case exact qualification of this hardware policy matched
raw SOC enumeration against its threshold formulas in all eight fixed cases,
including assignment ties and continuous response sets. All eighteen negative
controls and the smoke anchor passed. This is an internal implementation
qualification of the stated $n$-connector policy, not the fixed single-connector
model or an operational prevalence estimate.

For the main-text example in \eqref{eq:heteroexample}, the adequate-resource
policy also gives $U=15$ and the same cost-labelled projection. Its zero-early
endpoint has terminal demands 19, 15 and 15, fitting two 30 kW connectors by
serially grouping the last two demands. Its crossed full matching buys 9 early
on the 15-energy A bus and replenishes both buses by 20 terminally. These are
witnesses for this separate hardware policy. The main text supplies new serial
realizations at early power 15 and terminal power 49; the parallel witness
timings here are not reused as single-connector schedules. Individual allocation
sets can differ even when the joint load--cost projections coincide.
